% !TeX root = ../main.tex
\section{Badania symulacyjne}
\label{sec:simulation_results}

% ------------------------------------------------------------------------
% 5.1 Środowisko symulacyjne MATLAB i architektura eksperymentu
%     - struktura modułowa (step1_transmitter -> step2_channel -> step3_receiver -> step4_interpreter)
%     - centralna parametryzacja w get_wifi_params.m
%     - model kanału syntetycznego vs znormalizowany model wewnątrzbudynkowy TGax Model-B
%     - matryca 4 scenariuszy testowych (802.11a / 802.11ax x Ideal / TGax)
% 5.2 Porównanie charakterystyk detekcji: standard IEEE 802.11a a IEEE 802.11ax
%     - rozdzielczość odległościowa vs prędkość jednoznaczna v_unamb
%     - zjawisko aliasingu Dopplera w 802.11ax (dłuższy symbol OFDM)
% 5.3 Wpływ realistycznego kanału wielodrogowego TGax Model-B na jakość detekcji
%     - tłumienie stacjonarnego clutteru wewnątrzbudynkowego i przesłuchu bezpośredniego (MTI)
% 5.4 Efektywność detekcji w funkcji wejściowego stosunku sygnału do szumu (krzywe SNR)
% 5.5 Skuteczność algorytmu Coherent CLEAN w scenariuszach wielocelowych
% ------------------------------------------------------------------------

\subsection{Środowisko symulacyjne MATLAB i architektura eksperymentu}
\label{subsec:sim_environment}

\subsection{Porównanie charakterystyk detekcji: standard IEEE 802.11a a IEEE 802.11ax}
\label{subsec:sim_standards_comparison}

\subsection{Wpływ realistycznego kanału wielodrogowego TGax Model-B na jakość detekcji}
\label{subsec:sim_tgax_channel_effects}

\subsection{Efektywność detekcji w funkcji wejściowego stosunku sygnału do szumu}
\label{subsec:sim_snr_curves}

\subsection{Skuteczność algorytmu Coherent CLEAN w scenariuszach wielocelowych}
\label{subsec:sim_clean_performance}
